\documentclass[11pt]{article}

\usepackage{amsmath}
\usepackage{amssymb}
\usepackage{empheq}
\usepackage{geometry}
\usepackage{booktabs}
\usepackage{graphicx}

\geometry{margin=1in}

\usepackage[backend=biber,style=numeric,sorting=none]{biblatex}
\addbibresource{vector_potential.bib}

\newcommand{\Cbar}{{\bar C}}
\newcommand{\dd}{\mathrm{d}}
\newcommand{\pd}[2]{\frac{\partial #1}{\partial #2}}
\newcommand{\hx}{(1+hx)}
\newcommand{\bfA}{\mathbf{A}}
\newcommand{\eq}[1]{{(\protect\ref{#1})}}
\newcommand{\Eq}[1]{{Eq.~(\protect\ref{#1})}}
\newcommand{\Eqs}[1]{{Eqs.~(\protect\ref{#1})}}

\title{Vector Potential for the Generalized-Gradient Field Expansion\\
       in Frenet--Serret Coordinates with Finite Curvature}
\author{}
\date{}

\begin{document}
\maketitle

%---------------------------------------------------------------------------------------------------
\section{Introduction}

The work of Van der Schueren \textit{et al.}~\cite{tups09} generalizes the generalized
gradient formalism of Venturini and Dragt~\cite{venturini1999} to a not necessarily straight
curvilinear coordinate system. This paper describes how the work of Van der Schueren
can be applied to do symplectic tracking through bends.

%---------------------------------------------------------------------------------------------------
\section{Coordinates and Magnetic Field}

We work in the Frenet--Serret curvilinear coordinate system $(x,y,s)$ of
Van der Schueren~\cite{tups09}. To simplify the analysis, it is assumed that the reference
trajectory has constant curvature $h=1/\rho$ in the $x$--$s$ plane. This assumption is
appropriate for virtually all realistic simulations in high energy storage rings.

Following Van der Schueren, the magnetic field in the curvilinear coordinate
system is characterized by a set of functions $a_n(s)$, $b_n(s)$, and $b_s(s)$ where $n$ ranges
from one to infinity. Essentially, $n = 1$ is associated with dipole fields, $n = 2$ is associated
with quadrupole fields, etc. $b_s$ is associated with solenoid like fields. The association is
exact in the limit of $h = 0$ and the functions being independent of $s$.

The equations of Van der Schueren can be used to construct coefficients that characterize the
magnetic field:
\begin{align}
  B_x(x,y,s) = \sum_{i,j = 0,0}^{\infty,\infty} C_{x,i,j}(s) x^i y^j \label{bxc} \\
  B_y(x,y,s) = \sum_{i,j = 0,0}^{\infty,\infty} C_{y,i,j}(s) x^i y^j \label{byc} \\
  B_s(x,y,s) = \sum_{i,j = 0,0}^{\infty,\infty} C_{s,i,j}(s) x^i y^j . \label{bsc}
\end{align}
The coefficients $C$ are linear combinations of $a_n$, $b_n$, and $b_s$ and their derivatives.
Table xxx, which essentially reproduces Table 1 of Van der Schueren, show the lowest order coefficients.
Coefficients can easily be calculated using a symbolic math program and stored for use in tracking.

\begin{table}[h]
\centering
\caption{Lowest order $C_{x,i,j}$, $C_{y,i,j}$, and $C_{s,i,j}$ coefficients}
\renewcommand{\arraystretch}{2.4}
\setlength{\tabcolsep}{6pt}
\resizebox{\textwidth}{!}{%
\begin{tabular}{c c c c c c c}
\toprule
$\mathbf{(i,j)}$
   & $\mathbf{(0,0)}$ & $\boldsymbol{(1,0)}$ & $\boldsymbol{(0,1)}$
   & $\boldsymbol{(2,0)}$ & $\boldsymbol{(1,1)}$ & $\boldsymbol{(0,2)}$ \\
\midrule
$C_{x,i,j}$
 & $a_1$ & $a_2$ & $b_2$
 & $\dfrac{a_3}{2}$
 & $b_3$
 & $\dfrac{-a_3 - h\left(a_2 - h\,a_1 - 2b_s'\right) - a_1''}{2}$ \\
$C_{y,i,j}$
 & $b_1$ & $b_2$ & $-b_s' - h\,a_1 - a_2$
 & $\dfrac{b_3}{2}$
 & $-a_3 - h\left(a_2 - h\,a_1 - 2b_s'\right) - a_1''$
 & $-\dfrac{b_3 + h\,b_2 + b_1''}{2}$ \\
$C_{z,i,j}$
 & $b_s$ & $a_1' - h\,b_s$ & $b_1'$
 & $h^2 b_s - h\,a_1' + \dfrac{a_2'}{2}$
 & $b_2' - h\,b_1'$
 & $-\dfrac{a_2' + h\,a_1' + b_s''}{2}$ \\
\bottomrule
\end{tabular}%
}
\end{table}

%---------------------------------------------------------------------------------------------------
\section{Vector Potential}

For symplectic tracking the vector potential $\bfA$ is needed. In curvilinear coordinates
the the relationship between the vector potential and the field is
\begin{align}
  B_x = (\nabla\times\bfA)_x &= \pd{A_s}{y}
       - \frac{1}{1+hx}\,\pd{A_y}{s} \label{bxa} \\
  B_y = (\nabla\times\bfA)_y &= \frac{1}{1+hx}\,\pd{A_x}{s}
       - \frac{1}{1+hx}\,\pd{}{x}\!\bigl(\hx\,A_s\bigr) \label{bya}\\
  B_s = (\nabla\times\bfA)_s &= \pd{A_y}{x} - \pd{A_x}{y}. \label{bsa}
\end{align}
The vector potential is calculated by integrating the above equations along with \Eq{bxc} through
\Eq{bsc}. Evaluation at the integrals lower bounds 
 
Van der Schueren Eqs.~(8) through (10) is used to for the evaluation at the integrals lower bounds: 
\begin{align}
  B_x(x,y=0,s)   &= \sum_{n=1}^{\infty} a_n(s)\,\frac{x^{n-1}}{(n-1)!} \\[6pt]
  B_y(x,y=0,s)   &= \sum_{n=1}^{\infty} b_n(s)\,\frac{x^{n-1}}{(n-1)!} \\[6pt]
  B_s(x=0,y=0,s) &= b_s(s).
\end{align}

It is convenient to divide the calculation of the vector potential into three pieces by writing
\begin{align}
  C_{x,i,j}(s) &= \alpha_{x,i,j}(s) + \beta_{x,i,j}(s) + \gamma_{x,i,j}(s) \\
  C_{y,i,j}(s) &= \alpha_{y,i,j}(s) + \beta_{y,i,j}(s) + \gamma_{y,i,j}(s) \\
  C_{s,i,j}(s) &= \alpha_{s,i,j}(s) + \beta_{s,i,j}(s) + \gamma_{s,i,j}(s)
\end{align}
where the $\alpha$ terms are the part of $C$ that depend upon $a_n$, the $\beta$ terms
are the part of $C$ that depend upon $b_n$ and the $\gamma$ terms are the part of $C$
that depend upon $b_s$. For the $\alpha$ piece, the gauge where $A_y = 0$ is used to produce:
\begin{align}
  A_x &= -\sum_{i,j = 0,0}^{\infty,\infty} \frac{1}{j+1} \, \alpha_{s,i,j} \, x^{i} y^{j+1} \\
  A_y &= 0 \\
  A_s &=  \sum_{i,j = 0,0}^{\infty,\infty} \frac{1}{j+1} \, \alpha_{x,i,j} \, x^{i} y^{j+1}
\end{align}

For the $\beta$ part the same gauge $A_y = 0$ is used. Unlike the $\alpha$ piece, the $\beta$
piece has a nonzero value on the midplane ($B_y(x,0,s) = \sum_n b_n\,x^{n-1}/(n-1)!$), so an
additional $y$-independent term, built from the $\beta_{y,i,0}$ coefficients is needed
in $A_s$:
\begin{align}
  A_x &= -\sum_{i,j = 0,0}^{\infty,\infty} \frac{1}{j+1} \, \beta_{s,i,j} \, x^{i} y^{j+1} \\
  A_y &= 0 \\
  A_s &=  \sum_{i,j = 0,0}^{\infty,\infty} \frac{1}{j+1} \, \beta_{x,i,j} \, x^{i} y^{j+1}
       - \frac{1}{1 + h \, x} \sum_{i = 0}^{\infty} \, \beta_{y,i,0} \, \left(
         \frac{x^{i+1}}{i+1} + \frac{h \, x^{i+2}}{i+2} \right)
\end{align}

For the $\gamma$ part, the gauge where $A_z = 0$ is used to produce
\begin{align}
  A_x &=  (1 + h \, x) \sum_{i,j = 0,0}^{\infty,\infty} \left[ \int ds \, \gamma_{y,i,j} \right] \, 
          x^{i} y^{j} \\
  A_y &= -(1 + h \, x) \sum_{i,j = 0,0}^{\infty,\infty} \left[ \int ds \, \gamma_{x,i,j} \right] \, 
          x^{i} y^j \\
  A_s &= 0
\end{align}
The integrals are trivial since $\gamma$ is characterized by $b_s$ and its derivatives.

%--------------------------------------------------------------------------------------------------- 
\printbibliography

\end{document}
